% =============================================================
%  P2P Vision Challenge — Final Report
%  Authors: Vanshika Agrawal, Chiranjeev Kumar
%
%  Compiles with vanilla pdflatex (no acl.sty needed).
%  Tested: pdflatex p2p_report.tex  (run twice for refs)
% =============================================================

\documentclass[11pt,a4paper]{article}

\usepackage[T1]{fontenc}
\usepackage[utf8]{inputenc}
\usepackage{times}
\usepackage{microtype}
\usepackage{graphicx}
\usepackage{booktabs}
\usepackage{amsmath}
\usepackage{amssymb}
\usepackage{float}
\usepackage[hidelinks]{hyperref}
\usepackage[margin=0.85in]{geometry}
\usepackage{titlesec}
\usepackage{multicol}
\usepackage{abstract}

\titlespacing*{\section}{0pt}{0.85em}{0.25em}
\titlespacing*{\subsection}{0pt}{0.55em}{0.18em}
\titlespacing*{\paragraph}{0pt}{0.45em}{0.5em}
\renewcommand{\abstractnamefont}{\normalfont\bfseries}
\renewcommand{\abstracttextfont}{\normalfont\small}
\setlength{\columnsep}{0.5cm}
\setlength{\parskip}{1.5pt}

\title{\vspace{-1.3cm}\Large\textbf{Parameter-Efficient Fine-Tuning of SmolVLM-500M\\
       for Science Multiple-Choice Visual Question Answering}}
\author{\normalsize Vanshika Agrawal (\texttt{va2652}) \quad \textbullet \quad Chiranjeev Kumar (\texttt{ck4137}) \\[2pt]
        \small Pixels to Predictions: DL Vision Challenge}
\date{}

\begin{document}
\maketitle

\begin{abstract}
\noindent\small
We describe our submission to the \emph{Pixels to Predictions: DL Vision Challenge}, a science multiple-choice visual question answering task derived from ScienceQA. The competition fixes the backbone to \texttt{SmolVLM-500M-Instruct} and caps trainable parameters at five million. We fine-tune with QLoRA (4-bit NF4 weights, bfloat16 compute) using a rank-8, all-modules LoRA adapter that targets attention and MLP projections, yielding 4{,}784{,}128 trainable parameters (0.94\% of the model). Three independently seeded models are trained for three epochs each on an H100~80GB GPU; for each seed we additionally select the better of the epoch-2 and epoch-3 adapters by validation accuracy. Per-seed log-likelihood scores on the answer-letter token are then averaged across seeds for the final prediction. Our ensemble achieves a public-leaderboard accuracy of \textbf{0.82092}, an improvement of \textbf{+0.008} absolute over the best single seed (0.81287) and \textbf{+0.026} over the worst contributing seed (0.79476).
\end{abstract}

\vspace{0.1em}

\begin{multicols}{2}

\section{Introduction}

The \emph{Pixels to Predictions} challenge evaluates a fine-tuned vision--language model (VLM) on multiple-choice science questions paired with an image. Each instance provides a question, two-to-five textual answer choices, an image, and (for many instances) auxiliary text fields named \texttt{hint} and \texttt{lecture}. The task is to predict the zero-indexed correct choice; the leaderboard metric is plain classification accuracy.

Three rules constrain the design space: the backbone must be \texttt{SmolVLM-500M-Instruct}~\cite{marafioti2025smolvlm}; trainable parameters added during fine-tuning must not exceed five million; and no external data may be used. Within these constraints, the practical levers reduce to (i) the parameter-efficient fine-tuning configuration; (ii) the prompt format; (iii) the inference scoring procedure that maps model outputs to discrete answer indices; and (iv) any post-training combination of multiple models. Our recipe addresses each of these and achieves a public-leaderboard accuracy of \textbf{0.82092}.

\section{Dataset}

\subsection{Training Data}

The provided training split contains 3{,}109 image--question instances spanning three top-level subjects: \emph{natural science} (2{,}264; 72.8\%), \emph{social science} (768; 24.7\%), and \emph{language science} (77; 2.5\%). Grade levels are distributed across grades 1--12, with grades 4--8 making up 86\% of the training set. Each instance is annotated with hierarchical metadata fields (\texttt{subject}, \texttt{topic}, \texttt{category}, \texttt{skill}). The number of answer choices ranges from 2 to 5, with the mode being 3 choices (49.9\%); the answer-position distribution is heavily imbalanced toward earlier indices (position~0: 36.2\%, position~1: 33.1\%, position~2: 23.7\%, position~3: 6.6\%, position~4: 0.5\%).

\subsection{Data Preprocessing}

We apply minimal preprocessing in line with the competition rule prohibiting external data. (i) Images are loaded with PIL, converted to RGB, and resized via bicubic interpolation so that the longest side equals 384 pixels. The result is then center-pasted onto a $384\times384$ white background, preserving aspect ratio by letterboxing. (ii) The \texttt{hint} field (present in 76.7\% of train rows) is truncated to 600 characters; the \texttt{lecture} field (85.8\% of train rows) is truncated to 800 characters. These limits were chosen empirically to keep the prompt within a comfortable token budget while losing only the longest tail of pedagogical text. (iii) The \texttt{choices} field, stored as a JSON-encoded list of strings, is parsed and rendered as a lettered list (A, B, C, \ldots) in the prompt. No data augmentation is applied at training time.

\subsection{Validation and Test Splits}

The validation split contains 1{,}048 instances and the test split contains 1{,}008. The choice-count distribution and metadata coverage closely match the training split (test: hint coverage 78.8\%, lecture coverage 81.8\%), supporting the assumption that train-time conditioning on hint/lecture transfers to the test distribution.

\subsection{Visual Content}

A manual inspection of representative images reveals four broad image categories whose visual reasoning demands differ markedly: (i) photographs of organisms and physical objects (e.g., snakes, frogs, parachutes); (ii) labeled US/regional maps requiring identification of highlighted states or regions; (iii) ecosystem and food-web diagrams with white directed arrows between named entities; and (iv) physics/chemistry schematics with numeric labels (solution flasks, magnet pairs, particle samples). The first category is straightforward classification; the latter three require reading small-text labels and tracing relationships among them. Image dimensions vary widely (observed range $162{\times}162$ to $760{\times}506$, with extreme cases such as $606{\times}87$ for some magnet-pair questions, aspect ratio 6.97). Our $384{\times}384$ letterbox preserves aspect ratio but inevitably wastes canvas: ultra-wide magnet schematics use only 14--27\% of the input pixels after padding, while square photographs use the full 100\%. Magnet-related questions account for 11.7\% of training (365 of 3{,}109; all three-choice), so this preprocessing inefficiency, while not fatal, plausibly affects a non-trivial slice of the data.

\section{Model Description}

\subsection{Base Model}

We use the model identifier \texttt{SmolVLM-500M-} \texttt{Instruct} loaded via \texttt{transformers}~4.57.6 with \texttt{AutoModelForVision2Seq}. The model couples a SigLIP-based vision encoder with a small autoregressive language model and applies pixel-shuffle visual-token compression. The associated \texttt{AutoProcessor} handles both tokenization and image preprocessing. The base model is loaded in 4-bit NF4 quantization~\cite{dettmers2023qlora} with double quantization enabled and bfloat16 as the compute dtype.

\subsection{QLoRA Configuration}

Our LoRA adapter targets all linear projections inside each transformer block of the language model: the attention projections (\texttt{q\_proj}, \texttt{k\_proj}, \texttt{v\_proj}, \texttt{o\_proj}) and the MLP projections (\texttt{gate\_proj}, \texttt{up\_proj}, \texttt{down\_proj}). We use rank $r=8$, scaling $\alpha=32$, dropout 0.05, and no bias adaptation. The choice to adapt all modules rather than attention only is motivated by recent guidance~\cite{schulman2025loraregret} that, for tasks requiring representational adjustment beyond attention re-weighting, MLP projections are at least as important as attention projections. The base model is wrapped with \texttt{peft.prepare\_model\_for\_kbit\_training} prior to adapter injection, and gradient checkpointing is enabled. This configuration yields exactly \textbf{4{,}784{,}128} trainable parameters (verified from the saved \texttt{adapter\_model.safetensors} files, 215{,}872 below the 5M cap), or 0.94\% of the total model. The MLP projections (\texttt{gate\_proj}, \texttt{up\_proj}, \texttt{down\_proj}) account for 56.5\% of these parameters; the attention projections account for the remaining 43.5\%. The named-module match additionally captures \texttt{q/k/v\_proj} inside the first 12 layers of the SigLIP vision encoder (contributing 442{,}368 of the 4.78M total), which is incidental but useful: it lets the adapter adjust visual feature projections in addition to the text reasoning, while still leaving meaningful headroom under the cap.

\subsection{Prompt Format}

Each prompt is constructed by concatenating, in order: any non-empty \texttt{hint} (truncated to 600 characters), any non-empty \texttt{lecture} (truncated to 800 characters), the question text, a lettered choice list, and the fixed instruction \emph{``Reply with the single letter of the correct choice.''} The full text is wrapped in the SmolVLM chat template via \texttt{processor.apply\_chat\_template}, with a single image placeholder inserted into the user turn. During training the gold answer letter is appended as the assistant turn; during inference only the user turn is provided and the model's first-token distribution is used directly.

\section{Experimentation}

\subsection{Training Setup}

We train with AdamW at learning rate $2\times10^{-4}$ (an order of magnitude above typical full fine-tuning, following the short-LoRA-run guidance of~\cite{schulman2025loraregret}), a cosine annealing schedule across the full run, and gradient norm clipping at 1.0. The per-device batch size is 2 with 4-step gradient accumulation, yielding an effective batch size of 8. Each seed is trained for 3 epochs over the full 3{,}109-instance training split. Adapter weights are saved at the end of each epoch, and a rolling intermediate checkpoint of the latest two save-steps is maintained on Drive for resumption. All training was performed on a single NVIDIA H100~80GB GPU; one epoch over the full training set takes approximately 19~minutes (measured directly from the timestamps of the saved per-epoch adapters), so a full three-epoch run completes in $\approx$57~min and the three-seed pipeline in under 3~hours.

\subsection{Inference Strategy}

At inference we score each candidate answer by the log-probability of its letter token at the position immediately following the assistant-turn header. Concretely, given the prompt $P$ for a question with choices indexed by letters $L_1, \dots, L_K$, we compute
\begin{equation}
s_k = \log P_\theta(L_k \mid P, I),
\end{equation}
where $I$ is the input image and $\theta$ are the LoRA-adapted model parameters. Because the SmolVLM tokenizer assigns different token ids to a letter depending on whether it is preceded by a space, we evaluate both forms (\texttt{`A'} and \texttt{`\textvisiblespace A'}) and retain the larger of the two log-probabilities. The predicted answer is $\hat{y} = \arg\max_k s_k$, restricted to the valid choice indices for that question. Scoring at the first generation position rather than via free-form generation eliminates parsing failure as a source of error and produces a continuous score amenable to ensembling.

\subsection{Ablations}

Within our compute budget, we performed three ablations of practical interest.

\paragraph{(A) Ensemble vs.\ single seed.} Table~\ref{tab:results} compares each individual-seed submission to the score-averaged ensemble. The ensemble outperforms the best individual seed (0.81287) by $+0.008$ absolute and improves over the lowest contributing seed (0.79476) by $+0.026$. Quantitatively, the ensemble overrules the seed-42 prediction on 53 test rows (5.3\%), seed-7 on 74 rows (7.3\%), and seed-17 on 82 rows (8.1\%). Score-level averaging is a strict generalization of plurality voting and recovers calibrated predictions even when one seed is uncertain.

\paragraph{(B) Seed sensitivity and generalization.} Validation and test accuracies for each contributing seed are summarized in Table~\ref{tab:seeds}. The val--test gap is positive for every seed (test exceeds val by 0.95--2.41 points), indicating that none of the seeds is overfitting at three epochs and that the validation split is, if anything, slightly harder than the test split (consistent with its marginally lower hint coverage of 77.9\% vs 78.8\%). Validation accuracies span 0.777--0.803 (a 2.6-point spread on identical hyperparameters), confirming that even at fixed configuration the seed alone produces meaningful variance. The Pearson correlation between val and test accuracies is 0.79, strong enough to make val a reliable model-selection signal but not perfect: the val-best seed (42) is also test-best, but the val-second seed (17) is test-third, illustrating why ensembling adds value over picking the val-best seed alone. We also note that an unrelated re-run of seed~42 by the other team member produced a leaderboard score of only 0.71629 (vs 0.81287 in our final ensemble), a $0.10$-point single-seed spread that further motivates averaging across seeds rather than relying on any one run.

\paragraph{(C) Epoch-2 vs.\ epoch-3 selection.} For each seed, we save adapters at the end of epochs 2 and 3 and evaluate both on the validation split before selecting the better one for test inference. This safeguards against overfitting on the third epoch at negligible compute cost. In our final ensemble runs, epoch-3 adapters were selected for all three seeds (verified directly from the checkpoint files), indicating that three epochs at the chosen learning rate did not yet harm validation generalization on this dataset.

\section{Results}

\begin{table}[H]
\centering
\small
\setlength{\tabcolsep}{4pt}
\begin{tabular}{lcc}
\toprule
\textbf{Submission} & \textbf{LB Acc.} & \textbf{$\Delta$ vs.\ ens.} \\
\midrule
Initial baseline (data only)        & 0.42857 & $-0.392$ \\
Seed 42 (re-run)                     & 0.71629 & $-0.105$ \\
Seed 17                              & 0.79476 & $-0.026$ \\
Seed 7                               & 0.80080 & $-0.020$ \\
Seed 42 (used in ensemble)           & 0.81287 & $-0.008$ \\
\textbf{3-seed score ensemble}      & \textbf{0.82092} & --- \\
\bottomrule
\end{tabular}
\caption{Public-leaderboard accuracy across all submissions. The ensemble averages per-seed log-likelihood scores on the answer-letter token over seeds $\{42, 7, 17\}$ at validation-selected epochs.}
\label{tab:results}
\end{table}

\begin{table}[H]
\centering
\small
\setlength{\tabcolsep}{4pt}
\begin{tabular}{lccc}
\toprule
\textbf{Seed} & \textbf{Val acc} & \textbf{Test acc} & \textbf{Gap} \\
\midrule
42 & 0.8025 & 0.81287 & $+0.0104$ \\
7  & 0.7767 & 0.80080 & $+0.0241$ \\
17 & 0.7853 & 0.79476 & $+0.0095$ \\
\midrule
mean & 0.7882 & 0.80281 & $+0.0146$ \\
ens. & ---    & 0.82092 & --- \\
\bottomrule
\end{tabular}
\caption{Per-seed validation vs.\ public-test accuracy. Test exceeds val for every seed, indicating no overfitting at three epochs. Pearson $r(\text{val},\text{test})=0.79$.}
\label{tab:seeds}
\end{table}

\paragraph{Score decomposition.} The 0.82092 ensemble accuracy corresponds to roughly 827 correct predictions out of 1{,}008 test instances. Given the heavy answer-position prior in the training distribution (position~0 alone covers 36.2\%), even a degenerate constant predictor would score around 0.36, so a substantial fraction of our gains comes from genuine visual--textual reasoning rather than from memorizing the prior. The 0.42857 ``Initial baseline'' submission was an early dry-run that confirmed the submission pipeline was wired correctly.

\paragraph{Error analysis.} A breakdown of inter-seed predictions across the 1{,}008 test instances reveals a clear three-way pattern: all three seeds agree on \textbf{831 instances (82.4\%)}, two of the three agree on \textbf{173 instances (17.2\%)}, and all three produce different answers on only \textbf{4 instances (0.4\%)}. This $\sim$17\% middle band is precisely where score-level ensembling adds value. Pairwise agreement is highest between seeds~42 and~7 (90.4\%), and lowest between seeds~7 and~17 (86.6\%). Disagreement is also strongly correlated with the number of answer choices: pairwise disagreement averages 6.9\% on 2-choice questions, 15.7\% on 3-choice, 7.2\% on 4-choice, and \textbf{35.1\% on 5-choice} questions. Five-choice questions are both the rarest (3.8\% of test) and the hardest, consistent with the intuition that more choices both lower the chance baseline and amplify per-seed initialization sensitivity.

\section{Conclusion}

We presented a parameter-efficient recipe for SmolVLM-500M on a science multiple-choice VQA task. Within the competition's hard constraints (fixed backbone, 5M-parameter cap, no external data), we apply QLoRA at rank~8 across all attention and MLP projections, train three independently seeded models with epoch-level validation selection, and combine them by averaging answer-letter log-probabilities. The resulting submission scores \textbf{0.82092} on the public leaderboard.

\section{Code and Reproducibility}

\paragraph{Software.} \texttt{transformers}~4.57.6, \texttt{peft}~0.18.1, \texttt{bitsandbytes}~0.48.1, \texttt{accelerate}~1.10.1, PyTorch~2.x with bfloat16 + 4-bit NF4. All seeds (\texttt{random}, NumPy, PyTorch CPU, all CUDA devices) are explicitly set at the start of every training run.

\paragraph{Hardware.} Single NVIDIA H100~80GB; 4-bit NF4 quantization with bfloat16 compute; gradient checkpointing enabled. One epoch over 3{,}109 instances at effective batch size 8 takes $\approx$19~minutes; a full three-epoch seed completes in $\approx$57~minutes.

\paragraph{Reproduction.} Hyperparameters in Table~\ref{tab:hparams}. The pipeline is one Colab notebook with three cells: (1) setup (Drive mount, install, model load); (2) per-seed train and inference (run three times with \texttt{SEED} edited; produces \texttt{submission\_seed\{N\}.csv} and \texttt{test\_scores\_seed\{N\}.npy}); (3) ensemble (loads the three score arrays from Drive, masks invalid choice positions to $-\infty$, takes element-wise mean, and produces the final submission). Drive checkpointing is atomic (write to \texttt{.tmp}, rename) with rolling retention of the last two intermediate steps for resumption after disconnection. Notebook, saved adapters, per-seed prediction CSVs, raw \texttt{.npy} score arrays, and validation-accuracy logs are available at \url{https://drive.google.com/drive/folders/1epS_JpchcPya155hKrMfxR1w6OKR1I_8}.

\begin{table}[H]
\centering
\small
\setlength{\tabcolsep}{4pt}
\begin{tabular}{ll}
\toprule
\textbf{Hyperparameter} & \textbf{Value} \\
\midrule
Base model              & SmolVLM-500M-Instruct \\
Quantization            & 4-bit NF4, double-quant \\
Compute dtype           & bfloat16 \\
LoRA rank $r$           & 8 \\
LoRA $\alpha$           & 32 \\
LoRA dropout            & 0.05 \\
LoRA bias               & none \\
Target modules          & q,k,v,o,gate,up,down \\
Trainable parameters    & 4{,}784{,}128 \\
Optimizer               & AdamW \\
Learning rate           & $2\times10^{-4}$ \\
LR schedule             & Cosine annealing \\
Per-device batch size   & 2 \\
Gradient accumulation   & 4 \\
Effective batch size    & 8 \\
Epochs                  & 3 (ep3 selected for all seeds) \\
Max grad norm           & 1.0 \\
Image                   & $384^2$ letterbox, bicubic \\
Hint / lecture truncation & 600 / 800 chars \\
Seeds (ensemble)        & $\{42, 7, 17\}$ \\
Hardware                & 1$\times$ NVIDIA H100 80GB \\
\bottomrule
\end{tabular}
\caption{Final submission hyperparameters.}
\label{tab:hparams}
\end{table}

\begin{thebibliography}{9}
\small

\bibitem{dettmers2023qlora}
T.~Dettmers, A.~Pagnoni, A.~Holtzman, and L.~Zettlemoyer.
QLoRA: Efficient finetuning of quantized LLMs.
\emph{NeurIPS}, 2023.

\bibitem{hu2022lora}
E.~J.~Hu et al.
LoRA: Low-rank adaptation of large language models.
\emph{ICLR}, 2022.

\bibitem{lu2022scienceqa}
P.~Lu et al.
Learn to explain: Multimodal reasoning via thought chains for science question answering.
\emph{NeurIPS}, 2022.

\bibitem{marafioti2025smolvlm}
A.~Marafioti et al.
SmolVLM: Redefining small and efficient multimodal models.
\emph{arXiv:2504.05299}, 2025.

\bibitem{robinson2023leveraging}
J.~Robinson and D.~Wingate.
Leveraging large language models for multiple choice question answering.
\emph{ICLR}, 2023.

\bibitem{schulman2025loraregret}
J.~Schulman et al.
LoRA without regret: Practical recipes for parameter-efficient fine-tuning.
Technical report, 2025.

\end{thebibliography}

\end{multicols}

\end{document}
